\section*{第 \ref{ch:g4:classifying-animals} 章——动物分类}

\begin{solution}{exo:g4:classifying-animals:1}
\omterm{def:g4:classifying-animals:vertebrate}{脊椎动物}有以\omterm{def:g4:classifying-animals:vertebrate}{脊柱}为核心构成的内\omterm{def:g3:bones-and-muscles:skeleton}{骨骼}（如
鳟鱼、猫）；\omterm{def:g4:classifying-animals:vertebrate}{无脊椎动物}没有\omterm{def:g4:classifying-animals:vertebrate}{脊柱}（如蜗牛、
蜜蜂、蚯蚓、水母）。
\end{solution}

\begin{solution}{exo:g4:classifying-animals:2}
\omterm{prop:g3:sorting-living-things:groups}{鱼类}：有鳞片、鳍，在水中呼吸。\omterm{prop:g3:sorting-living-things:groups}{两栖动物}：湿润裸露的
\omterm{def:g1:the-five-senses:sense}{皮肤}，卵产于水中，幼年为\omterm{ex:g2:animals-grow-up:frog}{蝌蚪}期。\omterm{prop:g4:classifying-animals:classes}{爬行动物}：干燥带鳞的
\omterm{def:g1:the-five-senses:sense}{皮肤}，多数在陆地上产带壳的卵。\omterm{prop:g3:sorting-living-things:groups}{鸟类}：有\omterm{def:g1:animals-around-us:covering}{羽毛}、喙，
产带壳的卵。\omterm{prop:g3:sorting-living-things:groups}{哺乳动物}：有\omterm{def:g1:animals-around-us:covering}{皮毛}或毛发，幼体以\omterm{def:g2:animals-grow-up:born}{乳汁}为食。
\end{solution}

\begin{solution}{exo:g4:classifying-animals:3}
沙丁鱼：\omterm{prop:g3:sorting-living-things:groups}{鱼类}。鳄鱼：\omterm{prop:g4:classifying-animals:classes}{爬行动物}。蝾螈：\omterm{prop:g3:sorting-living-things:groups}{两栖动物}。鸵鸟：
\omterm{prop:g3:sorting-living-things:groups}{鸟类}。鲸：\omterm{prop:g3:sorting-living-things:groups}{哺乳动物}。蚯蚓：\omterm{def:g4:classifying-animals:vertebrate}{无脊椎动物}——没有\omterm{def:g4:classifying-animals:vertebrate}{脊柱}，所以
根本不属于任何\omterm{def:g4:classifying-animals:vertebrate}{脊椎动物}纲。
\end{solution}

\begin{solution}{exo:g4:classifying-animals:4}
猫在\omterm{prop:g3:sorting-living-things:groups}{哺乳动物}之中，\omterm{prop:g3:sorting-living-things:groups}{哺乳动物}在\omterm{def:g4:classifying-animals:vertebrate}{脊椎动物}之中，
\omterm{def:g4:classifying-animals:vertebrate}{脊椎动物}在\omterm{def:g1:animals-around-us:animal}{动物}之中。
\end{solution}

\begin{solution}{exo:g4:classifying-animals:5}
\omterm{def:g1:the-five-senses:sense}{皮肤}——蜥蜴的干燥带鳞，蝾螈的湿润裸露——
以及卵：蜥蜴的卵带壳、产于陆地，
蝾螈的卵无壳、产于水中。
\end{solution}

\begin{solution}{exo:g4:classifying-animals:6}
\omterm{def:g4:classifying-animals:vertebrate}{无脊椎动物}。例如：\omterm{prop:g3:sorting-living-things:groups}{昆虫}、蜘蛛、蜗牛、
蠕虫（还有水母）。
\end{solution}

\begin{solution}{exo:g4:classifying-animals:7}
企鹅：有\omterm{def:g1:animals-around-us:covering}{羽毛}吗？有——是\omterm{prop:g3:sorting-living-things:groups}{鸟类}；在第一个问题处即可判定。
海豚：有\omterm{def:g1:animals-around-us:covering}{羽毛}吗？没有。有\omterm{def:g1:animals-around-us:covering}{皮毛}或\omterm{def:g2:animals-grow-up:born}{乳汁}吗？它用\omterm{def:g2:animals-grow-up:born}{乳汁}哺育幼崽——
是\omterm{prop:g3:sorting-living-things:groups}{哺乳动物}；到此为止。
\end{solution}

\begin{solution}{exo:g4:classifying-animals:8}
例如：\omterm{def:g1:animals-around-us:covering}{羽毛}、喙、翅膀、带壳的卵——以及
\omterm{def:g4:classifying-animals:vertebrate}{脊柱}，因为\omterm{prop:g3:sorting-living-things:groups}{鸟类}是\omterm{def:g4:classifying-animals:vertebrate}{脊椎动物}。
\end{solution}

\begin{solution}{exo:g4:classifying-animals:9}
问题在“幼体以\omterm{def:g2:animals-grow-up:born}{乳汁}为食”处即可停止：鲸用\omterm{def:g2:animals-grow-up:born}{乳汁}哺育
幼崽，所以它是\omterm{prop:g3:sorting-living-things:groups}{哺乳动物}，无论是否几乎无毛。“生活在
海里”是\omterm{def:g2:where-animals-live:habitat}{栖息地}，而不是身体特征——而且
\omterm{def:g4:classifying-animals:classification}{分类}依据的正是\omterm{def:g1:animals-around-us:animal}{动物}\emph{拥有}的特征，
因为习性和住处可以被不相干的\omterm{def:g1:animals-around-us:animal}{动物}共享：鲨鱼
和海豚共享的是住址，而非类别。
\end{solution}

\begin{solution}{exo:g4:classifying-animals:10}
\omterm{def:g4:classifying-animals:vertebrate}{脊柱}。蛇的\omterm{def:g4:classifying-animals:vertebrate}{脊柱}由数百块小骨头组成——
是\omterm{def:g4:classifying-animals:vertebrate}{脊椎动物}，\omterm{def:g1:the-five-senses:sense}{皮肤}干燥带鳞，与\omterm{prop:g4:classifying-animals:classes}{爬行动物}相同；蚯蚓
完全没有\omterm{def:g4:classifying-animals:vertebrate}{脊柱}——是\omterm{def:g4:classifying-animals:vertebrate}{无脊椎动物}。缺少
腿只是缺少的一种特征；
\omterm{def:g4:classifying-animals:vertebrate}{脊柱}这一问题应首先提出。
\end{solution}

\begin{solution}{exo:g4:classifying-animals:11}
\omterm{prop:g3:sorting-living-things:groups}{哺乳动物}胜出：\omterm{def:g1:animals-around-us:covering}{皮毛}和用于哺育幼体的\omterm{def:g2:animals-grow-up:born}{乳汁}是\omterm{prop:g3:sorting-living-things:groups}{哺乳动物}的特征组合，
鸭嘴兽两者兼有——产卵只是其中一个不合群的特征，
\omterm{def:g4:classifying-animals:classification}{分类}时要权衡整组特征，而不是
让单一特征说了算；在某一方面特殊的\omterm{def:g1:animals-around-us:animal}{动物}，
仍归入其大多数特征所指向的类群。（鸭嘴兽
的怪异程度足以使产卵的\omterm{prop:g3:sorting-living-things:groups}{哺乳动物}
自成一个小的亚群。）
\end{solution}
